\chapter*{Lời Mở Đầu}
\addcontentsline{toc}{chapter}{Lời Mở Đầu}
\markboth{Lời Mở Đầu}{Lời Mở Đầu}

\begin{truyen}{Lời Mở Đầu}{Opening Words}
\chuhoa{Q}{uyển này không bắt đầu từ những câu hỏi học sinh hỏi thầy, mà bắt đầu từ những điều thầy muốn ngồi xuống nói chậm với học sinh.}
\textit{(This volume does not begin with the questions students ask teachers, but with the things a teacher wants to sit down and say slowly to students.)}

Có những điều các em đã nghe nhiều lần: phải cố gắng, phải học, phải sống tử tế.
\textit{(There are things students have heard many times: work hard, study, live with decency.)}

Nhưng có khi các em chưa thật sự nghe được, vì chúng thường được nói quá nhanh, quá chung, hoặc bằng giọng làm người nghe thấy mình đang bị ép hơn là được hiểu.
\textit{(But sometimes students have not truly heard them, because such things are often said too quickly, too generally, or in a tone that feels more forceful than understanding.)}

Quyển này chọn một giọng khác. Không lên lớp quá cao. Không làm quá nhiều khẩu hiệu. Mỗi bài ngắn như một lời thầy nói lại một điều quen thuộc nhưng cố nói gần hơn, thật hơn, và có ích hơn cho đường học lẫn đường lớn lên của một học sinh.
\textit{(This book chooses a different voice. It does not lecture from too high above. It avoids empty slogans. Each short piece is like a teacher saying something familiar again, but closer, truer, and more useful for both study and growing up.)}

Nếu em đang đi học, thầy mong em đọc quyển này như đang ngồi nghe một người lớn muốn nói chuyện đàng hoàng với mình.
\textit{(If you are a student, I hope you read this as though sitting with an adult who wants to speak properly with you.)}

Nếu em đang dạy học, thầy mong quyển này giúp mình nhớ lại học sinh cần được nói bằng một giọng người như thế nào.
\textit{(If you are a teacher, I hope this helps us remember the kind of human voice students need to hear.)}
\end{truyen}

\begin{baihoc}
Có những điều đúng nhưng phải được nói đúng giọng thì mới thật sự chạm vào người nghe.
\textit{(Some things are true, but they must be said in the right voice before they can truly reach the listener.)}
\end{baihoc}

\begin{ghinhoanh}
\textbf{Short English to remember:}

Truth reaches deeper when it is spoken in a human voice.

\end{ghinhoanh}

\begin{tuhocnhanh}
1. Vì sao một lời đúng vẫn có thể không chạm được học sinh? \textit{Why can a true message still fail to reach a student?}

2. Học sinh cần nghe bằng kiểu giọng nào? \textit{What kind of voice do students need to hear?}

3. Nếu giữ một tinh thần cho cả quyển này, em muốn giữ tinh thần nào? \textit{If you keep one spirit from this whole book, what should it be?}
\end{tuhocnhanh}

\ngancach
